% !TeX program = pdflatex
% Compile twice with pdfLaTeX. The bibliography is included below.
\documentclass[12pt,a4paper]{article}
\usepackage[T1]{fontenc}
\usepackage{lmodern,microtype}
\usepackage[margin=1in]{geometry}
\usepackage{array,tabularx,titlesec,setspace}
\usepackage[table]{xcolor}
\usepackage[numbers]{natbib}
\usepackage[hidelinks]{hyperref}
\hypersetup{pdftitle={Manual, Scripted and Agent-Assisted ORCA Workflows},
  pdfauthor={CHEN GENG},pdflang={en}}

% Replace \blankfield with the relevant details when ready.
\newcommand{\blankfield}{\rule{3.5cm}{0.35pt}}
\newcommand{\StudentName}{CHEN GENG}
\newcommand{\StudentID}{\blankfield}
\newcommand{\Supervisor}{\blankfield}
\newcommand{\SubmissionDate}{\blankfield} % DD/MM/YYYY

\singlespacing
\setlength{\parindent}{0pt}
\setlength{\parskip}{0.42em}
\setlength{\emergencystretch}{1.5em}
\setlength{\tabcolsep}{5pt}
\renewcommand{\arraystretch}{1.16}
\newcolumntype{Y}{>{\raggedright\arraybackslash}X}
\titleformat{\section}{\fontsize{12}{14}\selectfont\bfseries}{\thesection.}{0.5em}{}
\titlespacing*{\section}{0pt}{1.05em}{0.35em}
\newcommand{\subhead}[1]{\par\addvspace{0.4em}\textbf{#1}\par\nobreak}
\newcommand{\unitE}{\ensuremath{\mathrm{kJ\,mol^{-1}}}}
\newcommand{\doi}[1]{\href{https://doi.org/#1}{doi:\,#1}}
\raggedbottom
\begin{document}

\textbf{BG6023 MSc Research I}\hfill CA1 Research Proposal\par
\vspace{0.25em}
{\fontsize{12}{15}\selectfont\bfseries
Project Title\quad Manual, Scripted and Agent-Assisted ORCA Workflows\par}
\begin{tabularx}{\linewidth}{|p{0.456\linewidth}|Y|}
\hline
\rowcolor[HTML]{F2F2F2}\textbf{Field} & \textbf{Details} \\ \hline
Student name & \StudentName \\ \hline
Student ID & \StudentID \\ \hline
Supervisor & \Supervisor \\ \hline
Date of submission & \SubmissionDate \\ \hline
\end{tabularx}

\section{Introduction and Background}
In molecular density functional theory (DFT), preparing an input, obtaining
a converged solution and interpreting the resulting property are separate
tasks. Jacobs and Pollice tested whether large language models (LLMs) could write
executable ORCA inputs~\cite{jacobs}. Q-Planner coupled LLM planning with
deterministic execution~\cite{qplanner}. Quntur was benchmarked against
Claude Code using ORCA skills and the same base model~\cite{quntur}.
For a chemist with an established protocol, the practical question is whether
either agent can match a checked manual result more economically than a fixed
script. Answering it calls for measuring hands-on time as well as correctness.

After a calculation, a chemist must decide whether to reuse the output or
spend resources on another geometry. CENSO screens conformers inexpensively
before refining their energies~\cite{censo}. Under a fixed DFT budget,
an early energy could help decide which candidate to evaluate next.

\section{Problem Statement}
In my preliminary Auto-ORCA trial, seven of eight ORCA jobs terminated
normally, but none of seven requests met all its requirements. A butane
scan stopped one endpoint short. In separate ethanol pilots, the electronic
energy decreased in implicit water at both basis sets. Inferring a standard
hydration free energy would require further thermodynamic information.
Six \emph{n}-hexane geometries spanned $12.61$~\unitE{} in electronic energy.
These cases motivate independent checks of delivery, evidence reuse and
candidate selection.

\section{Aim, Objectives and Research Questions}
\textbf{Aim.} To compare the scientific reliability and cost of four ORCA
workflows and test two decisions made by the project's dedicated Auto-ORCA agent.

\textbf{1. Benchmark delivery.} Which automated workflows reproduce the
independently checked manual result on six tasks, and what operator time and
compute does each of the four workflows require?

\textbf{2. Test evidence reuse.} Does checking the requested quantity and
conditions increase the fraction of eight question pairs answered correctly?

\textbf{3. Test conformer selection.} Across six molecules, does seeing
each DFT energy before the next choice increase the chance of finding the
lowest-energy candidate within four queries, or reduce the energy missed?

\clearpage
\section{Methodology}
\subhead{Study 1 \enspace Four workflows on the same six ORCA tasks}
I will freeze two held-out cases in each of three families. The first asks
for a signed gas-to-water electronic energy difference at fixed geometry
(B3LYP/def2-SVP with a conductor-like polarizable continuum model, CPCM).
The second asks for an optimised carbonyl
structure, a minimum check and a C=O stretching frequency assigned from
the normal-mode displacements. The third asks for six vertical singlet
excitation energies and the brightest root by length-form oscillator strength.
I will use PBE0/def2-SVP for the last two families and the Tamm--Dancoff
approximation for excitations. Each case specifies charge, multiplicity,
coordinates, method and deliverables.

\begin{tabularx}{\linewidth}{@{}>{\raggedright\arraybackslash}p{0.27\linewidth}Y@{}}
\hline
\textbf{Workflow} & \textbf{Who prepares, executes and interprets the calculation} \\ \hline
Manual ORCA & I will write the inputs, run ORCA and read the raw outputs. A second chemistry reader will check my inputs and readings before they become the hidden reference. \\
Fixed script & A frozen script with no LLM takes the case parameters, runs its predefined protocol and extracts the requested quantities. I will record the time to prepare the script separately. \\
Dedicated Auto-ORCA & The project's specialised agent receives the natural-language task and coordinates, plans the work, launches ORCA and interprets the output. \\
General agent with ORCA Skill & A general coding agent receives the same text and coordinates, plus a frozen ORCA Skill containing instructions and helper scripts. It plans, executes and interprets independently. \\ \hline
\end{tabularx}

The manual and scripted workflows run once per case. Each agent runs twice in
fresh workspaces, with order counterbalanced. The script is fixed before the
held-out cases are revealed. Automated workflows cannot read the manual inputs,
answer keys or other runs. I will match solver version, numerical settings,
cores and time limits across workflows. I will use a common base LLM for the
two agents if both support it. Any model or backend difference will be recorded
and the comparison interpreted at the level of the configured systems.

A reviewer will score masked automated submissions against the checked key.
A pass requires a fresh ORCA execution, the correct geometry and conditions,
convergence, the requested physical interpretation and evidence in the raw output.
Agreement limits are $0.1$~\unitE{} for the signed energy difference,
$2$~cm$^{-1}$ for the assigned stretch, and $0.002$~eV per excitation with the
correct brightest root. A minimum requires a converged geometry and no
imaginary internal harmonic modes there. I will report passes per case, failures and retries, along with
elapsed time, CPU core-hours, active operator minutes and model tokens. The manual
run supplies a timing baseline and a checked reference. A separate execution
check addresses the incomplete runs observed by Liu et al.\ in a materials
DFT benchmark~\cite{dftbench}.

\clearpage
\subhead{Study 2 \enspace What can the archived outputs establish?}
For eight held-out neutral singlets I will prepare gas and CPCM-water
B3LYP/def2-SVP energies at identical coordinates. Two molecules will use the
four checked manual outputs from Study~1. In four question pairs the direction
of subtraction reverses. In two, a question about the electronic energy
difference is followed in a separate session by one about the standard
hydration free energy. In the final two pairs I will withhold the water output
and ask for either the available gas energy or the solvent effect.

Both Auto-ORCA conditions receive identical read-only files and must give
units and source lines or explain what is missing. Only one prompts an explicit
check of quantity, geometry, method and solvent. I will counterbalance order and
run each question twice per condition in a fresh session (64 sessions).
A pair passes only if both answers agree with the independently reviewed key.
Tolerances are $0.1$~\unitE{} for differences and $10^{-5}$~hartree for total
energies. An unsupported quantity requires withholding the value and naming
the missing evidence. Per repeat and condition there are 12 answerable and
four unanswerable questions. I will count avoidable rerun requests, omissions
and unsupported numbers against their respective denominators.

\subhead{Study 3 \enspace Which conformer merits the next calculation?}
For six additional neutral molecules containing only C, H and O, with at
least three rotatable bonds and at most eight heavy atoms, I will embed up
to 200 candidates with a fixed seed, relax them with
MMFF94, and retain the first eight distinct geometries after symmetry-aware
heavy-atom alignment (0.35~\AA{} cutoff). MMFF94 and GFN2-xTB will rank the
same geometries. I will calculate and withhold all 48 gas-phase
$r^2$SCAN-3c single-point energies. Each lookup consumes one of four DFT
queries. Using fixed coordinates makes the penalty for a poor choice directly
measurable.

Two fixed rules will pick the four lowest MMFF94 candidates or favour distinct
predefined torsion families in that order. In the latter rule any remaining
places go to the next lowest candidates. Two Auto-ORCA conditions receive both
rankings and family labels, start with the MMFF94 minimum and make four choices
in four rounds. Both see their chosen IDs and remaining budget. Only one sees
each returned DFT energy. I will run each rule once per molecule and each agent
condition three times (36 sessions). AutoDFT's comparison of adaptive and
fixed plans in VASP illustrates why computational cost should accompany
success rates~\cite{autodft}.

A hit means finding a geometry within $0.01$~\unitE{} of the minimum in the
eight-candidate set. The energy missed is the best queried energy minus the
full-set minimum. Incomplete sessions count as misses, with their query count
and best available energy reported. In Studies~2 and~3 I will average repeats
within each molecule, report paired differences and show how the mean changes
when each molecule is left out in turn.

\clearpage
\section{Expected Outcomes and Significance}
The four-workflow comparison will show where a fixed script suffices and where
either agent can complete a task with less hands-on work. If the workflows
agree numerically, differences in interventions, retries and elapsed time may
still matter to a laboratory. I expect evidence checks to help most when a
question changes the physical quantity, and energy feedback to help most when
the two inexpensive rankings disagree. I will publish the reference inputs,
scoring rules, anonymised answers and selection sequences where licences allow.
Six benchmark cases and one manual operator support case-level conclusions.
The conformer study measures electronic-energy rankings within fixed candidate
sets. Equilibrium populations would require further sampling and free energies.

\section{Timeline}
\begin{tabularx}{\linewidth}{|>{\raggedright\arraybackslash}p{0.30\linewidth}|c|c|Y|}
\hline
\rowcolor[HTML]{F2F2F2}\textbf{Task} & \textbf{Start (week)} & \textbf{End (week)} & \textbf{Output} \\ \hline
Select cases and methods & 1 & 2 & Matched tasks and scoring protocol \\ \hline
Pilot and freeze & 3 & 4 & Feasibility checks and fixed workflows \\ \hline
Run four workflows & 5 & 6 & Checked references, outputs and cost logs \\ \hline
Test reuse and selection & 7 & 8 & Paired answers and energy tables \\ \hline
Audit and analyse & 9 & 10 & Per-case scores and error analysis \\ \hline
Write up and present & 11 & 12 & Research paper and presentation \\ \hline
\end{tabularx}

\section{Resources}
I will use ORCA~6.1.1, Python, RDKit, xTB, two agent environments and a
second chemistry reader. The benchmark will use one core per job and a
1~GiB memory limit, with serial scheduling and 192~MB MaxCore per process.
Each pass over the six tasks needs eight solver starts if optimisation and
frequencies run together. The six passes need 48 starts. The two later
studies need 16 solvent/gas and 48 conformer starts. Sharing four manual
outputs gives 108 planned starts. A proposed ceiling of 128, including pilots
and retries, at 15 minutes per start bounds solver time at 32 hours.
The 124 agent sessions (24, 64 and 36) each allow at most eight model calls.

By week 4 I will check that both agents can launch ORCA under the stated
limits, time a representative case, review model credit and confirm an
independent reader. I will freeze inputs, script, Skill, model identifiers,
reference criteria and settings before evaluation. Ineligible molecules may
be replaced from a prespecified list before unblinding. Failed submissions
stay in the denominator. Budget or service interruptions will be reported
with coverage by arm. Human time is my own operating time. Hidden answers
and credentials will stay outside agent workspaces.

\clearpage
\renewcommand{\bibsection}{\section{References}}
\setlength{\bibsep}{0.65em}
\begin{thebibliography}{9}
\raggedright
\bibitem{jacobs}
P. F. Jacobs and R. Pollice. Developing large language models for quantum
chemistry simulation input generation. \emph{Digital Discovery} 4 (2025),
762--775. \doi{10.1039/D4DD00366G}.

\bibitem{qplanner}
R. Zhang, W. Gao and Y. Kurashige. Q-Planner: A harness system for quantum
chemistry agents. \emph{ACS Omega} (2026). \doi{10.1021/acsomega.6c05744}.

\bibitem{quntur}
J. B. P\'erez-S\'anchez, Y. Zou, J. A. Campos-Gonzalez-Angulo, et al.
El Agente Quntur: A research collaborator agent for quantum chemistry.
\emph{Chem} (2026), 103259. \doi{10.1016/j.chempr.2026.103259}.

\bibitem{censo}
S. Grimme, F. Bohle, A. Hansen, P. Pracht, S. Spicher and M. Stahn.
Efficient quantum chemical calculation of structure ensembles and free
energies for nonrigid molecules. \emph{Journal of Physical Chemistry A}
125 (2021), 4039--4054. \doi{10.1021/acs.jpca.1c00971}.

\bibitem{dftbench}
S. Liu, J. Yang, W. Wan, J. Han, J. Li and T. Deng. DFT-Agent Bench:
Scaffold effects and execution provenance in autonomous DFT.
\emph{arXiv} (2026), \href{https://arxiv.org/abs/2611.10629v1}{2611.10629v1}
(preprint).

\bibitem{autodft}
P. Yang, Z. Zhang, Y. Li, et al. AutoDFT: A closed-loop multi-agent framework
for autonomous DFT calculations. \emph{arXiv} (2026),
\href{https://arxiv.org/abs/2605.26179v2}{2605.26179v2} (preprint).
\end{thebibliography}
\end{document}
